\subsection{04 августа. Д.р. Каратакия}
\textit{Метеоусловия: весь день солнечно. }

\begin{figure}[h!]
	\centering
	\includegraphics[angle=0, width=0.7\linewidth]{../pics/maps/04_map-1.jpg}
	\label{fig:04_map}
	
\end{figure}

4:40 Подъем и сбор лагеря. Дождя ночью не было, и в целом погода нам благоволила. Завтракаем и в 6:43 Выходим. С песнями поднимаемся по крутому подъему. В 7:55 мы доходим до выполаживания. 

8:35 бродим через р. Саватор и приваливаемся, дожидаемся группы Алексея Остапива, идущей в сторону пер. Саватор,  и прощаемся с ней. 

\begin{figure}[h!]
	\centering
	\includegraphics[angle=0, width=0.7\linewidth]{../pics/days/04_1.jpg}
	\label{fig:04_1}
	\caption{ Прощание с группой Алексея Остапива у брода через р. Саватор.}
\end{figure}



В 9:25 выдвигаемся вверх по склону и за 20 минут добираемся до хребта с прекрасным видом на д.р. Чон-Кызыл-Суу. Дальше идем в д.р. Каратакия, бродим и в 10:26 встаем на обед.
 
\begin{figure}[h!]
	\centering
	\includegraphics[angle=0, width=0.7\linewidth]{../pics/days/04_2.jpg}
	\label{fig:04_2}
	\caption{ Вид с хребта на д.р. Чон-Кызыл-Суу}
\end{figure}

\begin{figure}[h!]
	\centering
	\includegraphics[angle=0, width=0.7\linewidth]{../pics/days/04_3.jpg}
	\label{fig:04_3}
	\caption{ Заходим в д.р Каратакия.}
\end{figure}

12:10 выдвигаемся. Поднимаемся по левому берегу р. Каратакия, встречаем там несколько стад коров. Идем медленно, тк у некоторых участников появляются симптомы горняшки. 

К 14:00 доходим до места ночевки, ставим лагерь на холме недалеко от ручья. Посреди кочек было выбрано относительно вытоптанное место. Однако, прежде чем поставить палатки, пришлось основательно его расчистить от коровьих лепешек. 

\begin{figure}[h!]
	\centering
	\includegraphics[angle=0, width=0.7\linewidth]{../pics/days/04_5.jpg}
	\label{fig:04_5}
	\caption{ Место ночевки.}
\end{figure}
В 16:00 ужинаем и в 18:00 отбой.
 Примерно в 22:00 пришло большое стадо коров.  Коровы начали наступать на края палатки и даже кусать тент. По очереди Руковод и Штурман выходили для вежливого разговора с проходимцами. В конечном итоге угрозы затискать и заобнимать коров сработали,  и после четвертого разговора они больше не возвращались. Наутро обнаружилось, что один из колышков палатки был беспощадно погнут могучим коровьим копытом.


\begin{tabular}{|c|c|}
	\hline
	ЧХВ: & 3:53 \\
	\hline
	ГХВ: & 5:33 \\
	\hline
\end{tabular}

\clearpage